\documentclass[12pt,a4paper]{article}

\usepackage[utf8]{inputenc}
\usepackage[T1]{fontenc}
\usepackage[brazil]{babel}
\usepackage{geometry}
\usepackage{graphicx}
\usepackage{amsmath}
\usepackage{booktabs}
\usepackage{hyperref}
\usepackage{listings}
\usepackage{xcolor}
\usepackage{float}
\usepackage{enumitem}
\usepackage{titlesec}
\usepackage{parskip}
\usepackage{fancyhdr}
\usepackage{caption}
\usepackage{subcaption}

\geometry{
  top=2.5cm,
  bottom=2.5cm,
  left=3cm,
  right=2.5cm
}

\hypersetup{
  colorlinks=true,
  linkcolor=blue!70!black,
  urlcolor=blue!70!black,
  citecolor=blue!70!black
}

\lstdefinestyle{matlab}{
  language=Matlab,
  basicstyle=\ttfamily\small,
  keywordstyle=\color{blue}\bfseries,
  commentstyle=\color{green!60!black},
  stringstyle=\color{red!70!black},
  numberstyle=\tiny\color{gray},
  numbers=left,
  stepnumber=1,
  frame=single,
  breaklines=true,
  captionpos=b
}

\lstdefinestyle{python}{
  language=Python,
  basicstyle=\ttfamily\small,
  keywordstyle=\color{blue}\bfseries,
  commentstyle=\color{green!60!black},
  stringstyle=\color{red!70!black},
  numberstyle=\tiny\color{gray},
  numbers=left,
  stepnumber=1,
  frame=single,
  breaklines=true,
  captionpos=b
}

\pagestyle{fancy}
\fancyhf{}
\fancyhead[L]{\small Previs\~ao da Vida \'{U}til de Baterias}
\fancyhead[R]{\small\thepage}
\renewcommand{\headrulewidth}{0.4pt}

\title{
  \vspace{-1cm}
  {\large Relat\'{o}rio T\'{e}cnico} \\[0.5em]
  \textbf{Previs\~ao Orientada a Dados da Vida \'{U}til de Baterias\\
  Antes da Degrada\c{c}\~ao de Capacidade} \\[0.5em]
  {\normalsize Reposit\'{o}rio: \texttt{data-driven-prediction-of-battery-cycle-life-before-capacity-degradation}}
}
\author{Pedro Lievra}
\date{Junho de 2026}

\begin{document}

\maketitle
\thispagestyle{fancy}

\tableofcontents
\newpage

%--------------------------------------------------------------------------
\section{Introdu\c{c}\~ao}

Este relat\'{o}rio documenta o trabalho realizado no reposit\'{o}rio
\texttt{data-driven-prediction-of-battery-cycle-life-before-capacity-degradation},
que disponibiliza o c\'{o}digo de processamento de dados associado ao artigo
cient\'{i}fico publicado na \emph{Nature Energy} (2019):

\begin{quote}
\textit{K.A. Severson, P.M. Attia et al., ``Data-Driven Prediction of Battery
Cycle Life Before Capacity Degradation'', Nature Energy, 2019.}
\end{quote}

O objetivo central do trabalho \'{e} demonstrar como prever a vida \'{u}til de baterias de
\'{i}ons de l\'{i}tio --- medida pelo n\'{u}mero de ciclos de carga/descarga at\'{e} a degrada\c{c}\~ao
cr\'{i}tica da capacidade --- utilizando apenas dados coletados nos primeiros ciclos de
opera\c{c}\~ao. A abordagem \'{e} totalmente orientada a dados (\textit{data-driven}), n\~ao
dependendo de modelos f\'{i}sico-qu\'{i}micos expl\'{i}citos.

\textbf{Nota:} O c\'{o}digo de modelagem preditiva est\'{a} sujeito a licen\c{c}a acad\^{e}mica e
dispon\'{i}vel mediante contato com Richard Braatz (\texttt{braatz@mit.edu}). Este
repositório publica apenas o código de processamento dos dados.

%--------------------------------------------------------------------------
\section{Conjunto de Dados}

\subsection{Origem e Disponibilidade}

Os dados foram coletados no MIT e est\~ao publicamente dispon\'{i}veis em
\url{https://data.matr.io/1/}. O conjunto cont\'{e}m medi\c{c}\~oes de baterias cil\'{i}ndricas
de \'{i}ons de l\'{i}tio (LFP/grafite, A123 Systems APR18650M1A, 1{,}1\,Ah nominais)
submetidas a diferentes pol\'{i}ticas de carga r\'{a}pida.

\subsection{Estrutura em Tr\^{e}s Lotes (\textit{Batches})}

Os dados est\~ao divididos em tr\^{e}s lotes coletados em per\'{i}odos distintos:

\begin{table}[H]
\centering
\caption{Resumo dos tr\^{e}s lotes de dados}
\begin{tabular}{@{}lllp{6cm}@{}}
\toprule
\textbf{Lote} & \textbf{Arquivo} & \textbf{Data} & \textbf{Descri\c{c}\~ao} \\
\midrule
Batch 1 & \texttt{2017-05-12\_batchdata\_...} & Mai/2017 &
  Primeiro conjunto de baterias; 5 c\'{e}lulas continuam no Batch~2 \\
Batch 2 & \texttt{2017-06-30\_batchdata\_...} & Jun/2017 &
  Continua\c{c}\~ao do Batch~1 para algumas c\'{e}lulas; c\'{e}lulas adicionais \\
Batch 3 & \texttt{2018-04-12\_batchdata\_...} & Abr/2018 &
  Terceiro conjunto; inclui c\'{e}lulas com pol\'{i}ticas de carga vari\'{a}vel \\
\bottomrule
\end{tabular}
\end{table}

\subsection{Categorias de Dados por C\'{e}lula}

Cada bateria (\textit{cell}) possui dados agrupados em tr\^{e}s categorias:

\begin{description}[leftmargin=2cm, style=nextline]
  \item[Descritores (\textit{Descriptors})]
    Pol\'{i}tica de carregamento (\textit{charging policy}), vida \'{u}til em ciclos
    (\textit{cycle life}), c\'{o}digo de barras (\textit{barcode}) e canal de teste
    (\textit{channel}).

  \item[Dados de Resumo (\textit{Summary})]
    Informa\c{c}\~oes por ciclo: n\'{u}mero do ciclo, capacidade de descarga
    ($Q_\text{discharge}$), capacidade de carga ($Q_\text{charge}$), resist\^{e}ncia
    interna (IR), temperatura m\'{a}xima ($T_\text{max}$), temperatura m\'{e}dia
    ($T_\text{avg}$), temperatura m\'{i}nima ($T_\text{min}$) e tempo de carregamento
    (\textit{chargetime}).

  \item[Dados de Ciclo (\textit{Cycle})]
    Medi\c{c}\~oes dentro de cada ciclo: tempo, capacidade de carga/descarga, corrente,
    tens\~ao, temperatura. Inclui tamb\'{e}m vetores derivados:
    \begin{itemize}
      \item \textbf{discharge\_dQ/dV}: diferencial de capacidade (dQ/dV) durante
        a descarga;
      \item \textbf{Qdlin}: capacidade de descarga interpolada linearmente em
        tens\~ao;
      \item \textbf{Tdlin}: temperatura interpolada linearmente em tens\~ao.
    \end{itemize}
\end{description}

%--------------------------------------------------------------------------
\section{Implementa\c{c}\~ao}

O reposit\'{o}rio fornece implementa\c{c}\~oes em duas linguagens: MATLAB (implementa\c{c}\~ao
original da an\'{a}lise) e Python (porta para acessibilidade ampliada).

\subsection{Arquivos do Reposit\'{o}rio}

\begin{table}[H]
\centering
\caption{Arquivos presentes no reposit\'{o}rio}
\begin{tabular}{@{}lp{9cm}@{}}
\toprule
\textbf{Arquivo} & \textbf{Descri\c{c}\~ao} \\
\midrule
\texttt{LoadData.m}         & Script MATLAB principal de carregamento e
                               pr\'{e}-processamento dos dados \\
\texttt{EC\_dQdV.m}         & Fun\c{c}\~ao MATLAB para c\'{a}lculo diferencial de capacidade
                               (dQ/dV) \\
\texttt{EC\_dVdQ.m}         & Fun\c{c}\~ao MATLAB para c\'{a}lculo diferencial de tens\~ao
                               (dV/dQ) \\
\texttt{Figure4\_mac.m}     & Script MATLAB para gera\c{c}\~ao das figuras diagn\'{o}sticas
                               (Figura 4/5 do artigo) \\
\texttt{BuildPkl\_Batch1.ipynb} & Notebook Python para constru\c{c}\~ao do arquivo
                               \texttt{.pkl} do Batch~1 \\
\texttt{BuildPkl\_Batch2.ipynb} & Notebook Python para constru\c{c}\~ao do arquivo
                               \texttt{.pkl} do Batch~2 \\
\texttt{BuildPkl\_Batch3.ipynb} & Notebook Python para constru\c{c}\~ao do arquivo
                               \texttt{.pkl} do Batch~3 \\
\texttt{Load Data.ipynb}    & Notebook Python demonstrando o carregamento dos
                               dados \texttt{.pkl} \\
\texttt{requirements.txt}   & Depend\^{e}ncias Python do projeto \\
\bottomrule
\end{tabular}
\end{table}

\subsection{Depend\^{e}ncias Python}

\begin{table}[H]
\centering
\caption{Depend\^{e}ncias do ambiente Python}
\begin{tabular}{@{}ll@{}}
\toprule
\textbf{Biblioteca} & \textbf{Vers\~ao} \\
\midrule
Python   & 3.7.10 \\
NumPy    & 1.22.0 \\
Pandas   & 1.2.3  \\
SciPy    & 1.6.1  \\
Matplotlib & 3.3.4 \\
h5py     & 2.8.0  \\
Jupyter  & (\'{u}ltima dispon\'{i}vel) \\
\bottomrule
\end{tabular}
\end{table}

Para criar o ambiente com Conda:
\begin{lstlisting}[language=bash, numbers=none, caption={Cria\c{c}\~ao do ambiente conda}]
conda create --name battery_env --file requirements.txt
\end{lstlisting}

%--------------------------------------------------------------------------
\section{Pipeline de Processamento dos Dados}

\subsection{Carregamento e Mesclagem dos Lotes (MATLAB)}

O script \texttt{LoadData.m} executa as etapas a seguir.

\subsubsection{Carregamento do Batch 1}

\begin{lstlisting}[style=matlab, caption={Carregamento do primeiro lote}]
load('.\Data\2017-05-12_batchdata_updated_struct_errorcorrect')
batch1 = batch;
numBat1 = size(batch1,2);
\end{lstlisting}

\subsubsection{Mesclagem de C\'{e}lulas Continuadas do Batch 2 ao Batch 1}

Cinco c\'{e}lulas do Batch~1 continuaram sendo cicladas no Batch~2. Seus dados
s\~ao concatenados antes do processamento:

\begin{lstlisting}[style=matlab, caption={Concatena\c{c}\~ao de ciclos entre lotes}]
add_len = [661, 980, 1059, 207, 481];
batch2_idx = [8:10, 16:17];
for i = 1:5
    batch1(i).cycles(end+1:end+add_len(i)+1) = ...
        batch(batch2_idx(i)).cycles;
    % ... (campos de summary tambem concatenados)
end
\end{lstlisting}

\subsubsection{Filtragem e Limpeza do Batch 3}

O Batch~3 requer etapas adicionais de limpeza de dados:

\begin{enumerate}
  \item Remo\c{c}\~ao do canal~46 (canal~38 no array 0-indexado), que apresentou
    problemas na coleta;
  \item Remo\c{c}\~ao de c\'{e}lulas que n\~ao completaram a degrada\c{c}\~ao abaixo de
    $Q < 0{,}885$\,Ah ao final do teste;
  \item Remo\c{c}\~ao das baterias ruidosas do Batch~8 (\'{i}ndices 3, 40 e 41).
\end{enumerate}

\begin{lstlisting}[style=matlab, caption={Filtragem do Batch 3}]
batch3(38) = [];   % canal 46 com problema de coleta
endcap3 = zeros(numBat3,1);
for i = 1:numBat3
    endcap3(i) = batch3(i).summary.QDischarge(end);
end
rind = find(endcap3 > 0.885);
batch3(rind) = [];

nind = [3, 40:41]; % baterias ruidosas
batch3(nind) = [];
\end{lstlisting}

\subsubsection{Combina\c{c}\~ao Final}

Os tr\^{e}s lotes s\~ao combinados em um \'{u}nico array e c\'{e}lulas que n\~ao finalizaram
os testes no Batch~1 s\~ao opcionalmente removidas:

\begin{lstlisting}[style=matlab, caption={Combina\c{c}\~ao dos lotes}]
batch_combined = [batch1, batch2, batch3];
batch_combined([9,11,13,14,23]) = []; % celulas incompletas do Batch 1
\end{lstlisting}

\subsection{Vari\'{a}vel de Sa\'{i}da --- Vida \'{U}til em Ciclos}

A vari\'{a}vel de sa\'{i}da do modelo \'{e} o n\'{u}mero de ciclos at\'{e} a capacidade de descarga
atingir o limiar de 0{,}88\,Ah (aproximadamente 80\% da capacidade nominal):

\begin{lstlisting}[style=matlab, caption={Extra\c{c}\~ao da vari\'{a}vel alvo}]
bat_label = zeros(numBat,1);
for i = 1:numBat
    if batch_combined(i).summary.QDischarge(end) < 0.88
        bat_label(i) = find( ...
            batch_combined(i).summary.QDischarge < 0.88, 1);
    else
        bat_label(i) = size(batch_combined(i).cycles,2) + 1;
    end
end
\end{lstlisting}

Esse crit\'{e}rio define a ``vida \'{u}til'' de forma objetiva e padronizada, permitindo
compara\c{c}\~oes diretas entre c\'{e}lulas com diferentes pol\'{i}ticas de carregamento.

\subsection{Divis\~ao Treino/Teste}

A divis\~ao segue o protocolo descrito no artigo original:

\begin{lstlisting}[style=matlab, caption={Divis\~ao treino/teste}]
test_ind  = [1:2:(numBat1+numBat2), 84];
train_ind = 1:(numBat1+numBat2);
train_ind(test_ind) = [];
secondary_test_ind = numBat-numBat3+1:numBat;
\end{lstlisting}

\begin{itemize}
  \item \textbf{Treino prim\'{a}rio}: c\'{e}lulas dos Batches 1 e 2 n\~ao selecionadas para
    o teste.
  \item \textbf{Teste prim\'{a}rio}: c\'{e}lulas alternadas dos Batches 1 e 2 (a cada 2)
    mais a c\'{e}lula de \'{i}ndice 84.
  \item \textbf{Teste secund\'{a}rio}: todas as c\'{e}lulas do Batch~3 (teste de
    generaliza\c{c}\~ao para dados coletados independentemente).
\end{itemize}

%--------------------------------------------------------------------------
\section{Engenharia de Atributos}

\subsection{Diferencial de Capacidade --- dQ/dV}

A fun\c{c}\~ao \texttt{EC\_dQdV.m} calcula o diferencial de capacidade em rela\c{c}\~ao \`{a}
tens\~ao durante a descarga. O c\'{a}lculo \'{e} baseado em diferen\c{c}as finitas progressivas
com passo fixo de tens\~ao:

\begin{equation}
\frac{dQ}{dV}\bigg|_{V_k} \approx \frac{\Delta Q}{\Delta V} =
\frac{Q(i+j) - Q(i)}{E_{we}(i+j) - E_{we}(i)}
\end{equation}

onde $j$ \'{e} o menor \'{i}ndice tal que $|E_{we}(i+j) - E_{we}(i)| \geq \delta_V$,
com $\delta_V = 4\,\text{mV}$ por padr\~ao. A capacidade \'{e} normalizada pela
capacidade final $Q_\text{scale}$ antes do c\'{a}lculo. Zeros e NaNs iniciais s\~ao
removidos.

\subsection{Diferencial de Tens\~ao --- dV/dQ}

A fun\c{c}\~ao \texttt{EC\_dVdQ.m} implementa o inverso: diferencial de tens\~ao em
rela\c{c}\~ao \`{a} capacidade, com passo padr\~ao de $\delta_Q = 1\%$ da capacidade
normalizada:

\begin{equation}
\frac{dV}{dQ}\bigg|_{Q_k} \approx \frac{\Delta V}{\Delta Q} =
\frac{E_{we}(i+j) - E_{we}(i)}{Q(i+j) - Q(i)}
\end{equation}

A capacidade \'{e} invertida (\texttt{flipud}) para que a descarga seja processada
no sentido crescente de $Q$.

\subsection{Interpreta\c{c}\~ao F\'{i}sica dos Atributos}

As curvas de dQ/dV e dV/dQ s\~ao sens\'{i}veis a transforma\c{c}\~oes de fase nos eletrodos
da bateria. Mudan\c{c}as na forma ou posi\c{c}\~ao dos picos dessas curvas ao longo dos
ciclos indicam degrada\c{c}\~ao interna (perda de material ativo, aumento de resist\^{e}ncia)
--- mesmo quando a curva de capacidade total ainda n\~ao mostra queda significativa.

A diferen\c{c}a entre as curvas dos ciclos 10 e 100, em particular
$\Delta Q(V) = Q_{100}(V) - Q_{10}(V)$, mostrou-se o preditor mais poderoso
identificado pelo estudo:

\begin{equation}
\Delta Q(V) = Q_{\text{ciclo 100}}(V) - Q_{\text{ciclo 10}}(V)
\end{equation}

A vari\^{a}ncia de $\Delta Q(V)$ ao longo do perfil de tens\~ao \'{e} o principal
atributo de entrada para os modelos de regress\~ao descritos no artigo.

%--------------------------------------------------------------------------
\section{Visualiza\c{c}\~ao Diagn\'{o}stica}

O script \texttt{Figure4\_mac.m} gera figuras diagn\'{o}sticas em quatro pain\'{e}is
para cada c\'{e}lula analisada:

\begin{enumerate}
  \item \textbf{Painel superior --- dQ/dV (C/10):} curvas de diferencial de
    capacidade obtidas com corrente lenta (C/10) nos ciclos 1, 100 e no fim de
    vida.
  \item \textbf{Painel 2 --- dV/dQ (C/10):} curvas de diferencial de tens\~ao com
    corrente lenta nos mesmos ciclos de refer\^{e}ncia.
  \item \textbf{Painel 3 --- dQ/dV (4C):} curvas de diferencial de capacidade com
    corrente r\'{a}pida (4C) usada no ciclamento acelerado.
  \item \textbf{Painel inferior --- $\Delta Q(V)$:} diferen\c{c}a acumulada de
    capacidade em fun\c{c}\~ao da tens\~ao ao longo dos ciclos, desde o ciclo 10 at\'{e}
    o fim de vida, colorida pela vida \'{u}til em escala logar\'{i}tmica.
\end{enumerate}

Tr\^{e}s exemplos de baterias s\~ao apresentados com pol\'{i}ticas de carga distintas:
4C, 6C e 8C, correspondendo a tempos de carga de 15\,min, 10\,min e 7{,}5\,min
respectivamente.

%--------------------------------------------------------------------------
\section{Formato dos Dados --- MATLAB vs.\ Python}

\begin{table}[H]
\centering
\caption{Equival\^{e}ncia entre formatos de arquivo}
\begin{tabular}{@{}lll@{}}
\toprule
\textbf{Aspecto} & \textbf{MATLAB (.mat)} & \textbf{Python (.pkl)} \\
\midrule
Estrutura principal & \texttt{struct}  & Dicion\'{a}rio aninhado \\
Acesso aos ciclos   & \texttt{batch(i).cycles(j).campo} &
  \texttt{batch[cell\_key]['cycles'][j]['campo']} \\
Dados de resumo     & \texttt{batch(i).summary.QDischarge} &
  \texttt{batch[cell\_key]['summary']['QDischarge']} \\
Constru\c{c}\~ao          & Arquivos fornecidos pelo dataset & Notebooks \texttt{BuildPkl\_*} \\
\bottomrule
\end{tabular}
\end{table}

Os notebooks \texttt{BuildPkl\_Batch1.ipynb}, \texttt{BuildPkl\_Batch2.ipynb} e
\texttt{BuildPkl\_Batch3.ipynb} leem os arquivos \texttt{.mat} via \texttt{h5py} e
geram os arquivos \texttt{.pkl} correspondentes, tornando os dados acess\'{i}veis
diretamente em Python sem depend\^{e}ncia do MATLAB.

%--------------------------------------------------------------------------
\section{Hist\'{o}rico de Contribui\c{c}\~oes}

O reposit\'{o}rio recebeu contribui\c{c}\~oes ao longo de aproximadamente quatro anos
(2019--2022), com os seguintes marcos principais:

\begin{table}[H]
\centering
\caption{Principais marcos do hist\'{o}rico de commits}
\begin{tabular}{@{}llp{8cm}@{}}
\toprule
\textbf{Data} & \textbf{Autor} & \textbf{Contribui\c{c}\~ao} \\
\midrule
Mar/2019 & K. Severson / R. Braatz &
  Commit inicial: ferramentas de processamento, notebooks de constru\c{c}\~ao dos pkl \\
Mar/2019 & K. Severson &
  Atualiza\c{c}\~ao para mesclar dados do Batch~2 no Batch~1 \\
Abr/2019 & M. Hwasser &
  Corre\c{c}\~ao de typo na extra\c{c}\~ao de dados \\
Jun/2020 & P. Attia &
  Adi\c{c}\~ao das fun\c{c}\~oes \texttt{EC\_dQdV.m} e \texttt{EC\_dVdQ.m}; script
  \texttt{Figure4\_mac.m} \\
Jan/2021 & P. Attia &
  Documenta\c{c}\~ao de \texttt{Qdlin} e \texttt{Tdlin} no README \\
Mai/2021 & P. Herring (TRI) &
  Corre\c{c}\~ao da l\'{o}gica de remo\c{c}\~ao de c\'{e}lulas (depend\^{e}ncia de ordem no MATLAB) \\
Mai/2021 & J. Schaeffer (ETH Z\"urich) &
  Cria\c{c}\~ao do \texttt{requirements.txt}; atualiza\c{c}\~ao do README \\
Nov/2022 & dependabot / P. Attia &
  Atualiza\c{c}\~ao de seguran\c{c}a: NumPy de 1.19.2 para 1.22.0 (CVE fix) \\
\bottomrule
\end{tabular}
\end{table}

%--------------------------------------------------------------------------
\section{Conclus\~oes}

O reposit\'{o}rio entrega uma base s\'{o}lida de processamento de dados para o problema
de previs\~ao de vida \'{u}til de baterias. Os pontos-chave do trabalho realizado s\~ao:

\begin{enumerate}
  \item \textbf{Coleta e organiza\c{c}\~ao} de dados de 124 baterias LFP/grafite em
    tr\^{e}s lotes independentes, com limpeza criteriosa de c\'{e}lulas defeituosas.

  \item \textbf{Unifica\c{c}\~ao} dos dados originalmente distribu\'{i}dos em dois lotes
    sequenciais, preservando a continuidade dos ciclos de c\'{e}lulas que
    cruzaram os per\'{i}odos de coleta.

  \item \textbf{Engenharia de atributos} centrada em curvas de dQ/dV e dV/dQ
    --- representa\c{c}\~oes fisicamente interpret\'{a}veis que capturam a degrada\c{c}\~ao
    interna antes que ela seja vis\'{i}vel na capacidade integrada.

  \item \textbf{Padroniza\c{c}\~ao} da vari\'{a}vel alvo: n\'{u}mero de ciclos at\'{e}
    $Q < 0{,}88$\,Ah, equivalente a aproximadamente 80\% da capacidade nominal.

  \item \textbf{Portabilidade}: os dados originalmente em MATLAB foram
    convertidos para Python (.pkl), tornando o trabalho acess\'{i}vel \`{a} comunidade
    sem licen\c{c}a de MATLAB.

  \item \textbf{Protocolo de valida\c{c}\~ao rigoroso}: a divis\~ao treino/teste \'{e}
    definida em c\'{o}digo, garantindo reprodutibilidade dos resultados do artigo
    publicado.
\end{enumerate}

O c\'{o}digo de modelagem --- que implementa regress\~ao regularizada sobre os atributos
descrito e atinge erros de previs\~ao inferiores a 10\% na vida \'{u}til de ciclos ---
est\'{a} dispon\'{i}vel mediante licen\c{c}a acad\^{e}mica junto aos autores do artigo.

%--------------------------------------------------------------------------
\section*{Refer\^{e}ncia}

K.A.~Severson, P.M.~Attia, N.~Jin, N.~Perkins, B.~Jiang, Z.~Yang,
M.H.~Chen, M.~Aykol, H.K.~Herring, D.~Fraggedakis, M.Z.~Bazant,
S.J.~Harris, W.C.~Chueh, R.D.~Braatz,
\textit{Data-driven prediction of battery cycle life before capacity degradation},
\textbf{Nature Energy} 4, 383--391 (2019).
\url{https://doi.org/10.1038/s41560-019-0356-8}

\end{document}
